\chapter{Lugares geométricos}
\begin{tikzpicture}
	\fill [left color=red!50, right color=teal!50] (0,0) rectangle (6.5,.2);
	\fill [left color=teal!50, right color=blue!50] (6.5,0) rectangle (11.5,.2);
	\end{tikzpicture}
\vspace{1cm}

\begin{adjustwidth}{40pt}{40pt}
\begin{cuadro-gris}

	\begin{multicols}{2}
	$\triangleright \quad$  Lugar geométrico.
	
		$\triangleright \quad$  Bisectriz.
	
	$\triangleright \quad$  Mediatriz.
	
	$\triangleright \quad$  Circunferencia.
	
	\end{multicols}
	
\end{cuadro-gris}
\end{adjustwidth}

\vspace{1cm}
\section{Lugares geométricos}
\begin{tikzpicture}
	\fill [left color=red!50, right color=teal!50] (0,0) rectangle (3.5,.1);
	\fill [left color=teal!50, right color=blue!50] (3.5,0) rectangle (7.5,.1);
	\end{tikzpicture}
\vspace{1cm}

	\begin{definition}[ Lugares geométricos]
	
	Conjunto de los puntos $X(x,y)$ del plano, $\mathbb R^2$	, que cumplen una determinada condición.
	\end{definition}

\vspace{5mm}
\underline{Observaciones}: $\quad$ Rectas y segmentos también son lugares geométricos.

\begin{itemize}
\item Una recta, $r:\, \begin{cases}\ A(x_0,y_0) \\ \ \vec v=(v_x,v_y) \end{cases} \ \to \ $

\hspace{2cm} $\to \quad r:\ X(x,y)\in \mathbb R^2 \ / \ X=A+\lambda \vec v\, , \ \ \forall t\in \mathbb R$

Si cambiamos la condición de igualdad por una desigualdad se obtiene un '`semiplano'.
\item Segmento $AB$, $\ r:\, \begin{cases} \ A \\ \ B \to \ \vec v =\overrightarrow{AB}=B-A \end{cases} \ \to \ $
 
\hspace{2cm} $\to \quad \overline{AB}:\ X\in \mathbb
 R^2 \ / \ X=A+\lambda \overrightarrow{AB} \ \text{ con } \ \lambda \in [0,1]$
 \end{itemize}

\vspace{5mm}
\begin{large}
\textbf{Tipos de lugares geométricos}	
\end{large}

\subsection{Mediatriz}
\begin{tikzpicture}
	\fill [left color=red!50, right color=teal!50] (0,0) rectangle (3.5,.01);
	\fill [left color=teal!50, right color=blue!50] (3.5,0) rectangle (7.5,.01);
	\end{tikzpicture}
\vspace{0.5cm}

\begin{definition}[ Mediatriz]

\emph{Mediatriz} de un segmento $AB$  es el lugar geométrico de los puntos del plano que equidistan de los puntos $A$ y $B$, extremos del segmento.	
\end{definition}

\vspace{5mm}

\color{gris}
\rule{250pt}{0.1pt}

$A(x_0,y_0), \ B(x_1,y_1)\quad $ Sea $X(x,y)\in \mathbb R^2$ un punto cualquiera del plano.

$X$ es un punto de la mediatriz de $AB \ \leftrightarrow \ d(X,A)=d(X.B) \quad \Rightarrow$

$\Rightarrow \ \sqrt{(x-x_0)^2+(y-y_0)^2} =  \sqrt{(x-x_1)^2+(y-y_1)^2} \ \to \ 
(x-x_0)^2+(y-y_0)^2=(x-x_1)^2+(y-y_1)^2 \ \to \
\cancel{x^2}-2xx_0+x_0^2+\bcancel{y^2}-2yy_0+y_0^2=\cancel{x^2}-2xx_1+x_1^2+\bcancel{y^2}-2yy_1+y_1^2 \ \to \ x(-2x_0+2x_1)+y(-2y_0+2y_1)=x_1^2+y_1^2-x_0^2-y_0^2 \ \ \Rightarrow \ \ $

\hspace{25mm}$\Rightarrow \quad  med:\  x(-2x_0+2x_1)+y(-2y_0+2y_1)-(x_1-x_0)(x_1+x_0)-(y_1-y_0)(y_1+y_0) $

Por otro lado, $med:\, \begin{cases} \ M=\dfrac{A+B}{2}=\left( \dfrac{x_1+x_0}{2},  \dfrac{y_1+y_0}{2} \right) \\ \ \perp \overrightarrow{AB} \ \to \ \vec v_{med}=\big( -(y_1+y_0),x_1-x_0) \end{cases}  \to \  \dfrac{x- \dfrac{x_1+x_0}{2}}{-(y_1-y_0)}=\dfrac{y- \dfrac{y_1+y_0}{2}}{x_1-x_0} \ \to \ \dfrac{2x-(x_0+x_1)}{-\cancel{2}(y_1-y_0)}=\dfrac{2y-(y_0+y_1)}{\cancel{2}(x_1-x_0)} \ \ \Rightarrow  \ \ $

\hspace{25mm} $\Rightarrow \quad  med:\  x(-2x_0+2x_1)+y(-2y_0+2y_1)-(x_1-x_0)(x_1+x_0)-(y_1-y_0)(y_1+y_0) $

\vspace{-5mm}
\begin{flushright} \rule{250pt}{0.1pt} \end{flushright}
\color{black}

\vspace{5mm}


\begin{multicols}{2}
Gráficamente, la mediatriz es una recta perpendicular a un segmento que lo divide en dos partes iguales.


Como podemos observar en la imagen, $med_{\, AB}$ divide al segmento $AB$ en dos partes iguales. 

     \begin{equation*} d(X,A)=d(X,B)\end{equation*} 
     
\begin{figure}[H]
	\centering
	\includegraphics[width=0.4\textwidth]{img-lg/lg01.png}
\end{figure}
\end{multicols}

Para el cálculo de una mediatriz, encontrar $\ med_{\, AB}:\, \begin{cases} \ M_{\, AB} \\ \perp \overrightarrow{AB} \end{cases}$ será más rápido que aplicar la definición, $\ med_{AB} = \{ X\in \mathbb R^2 \ \big/ \ d(X.A)=d(X,B)  \}$
 
\begin{miejemplo}
	


Encontrar la ecuación de la mediatriz del segmento con extremos en A(2 , 5) y B(4, -7).

\vspace{5mm} Para encontrar la ecuación de la recta de la mediatriz, usaremos la definición, es decir sabemos que la recta debe satisfacer que
$\quad d(A,X)=d(B,X) \quad \text{con } \ X(x,y) $
     
\vspace{2mm} Sustituyendo con los valores dados tenemos:$\quad \sqrt{(x-2)^2+(y-5)^2}=\sqrt{(x-4)^2+(y+7)^2}$
     
\vspace{2mm} Elevando al cuadrado en ambos lados de la ecuación y desarrollando los binomios tenemos: $\quad x^2-4x+4+y^2-10y+25=x^2-8x+16+y^2+14y+49$ 

\vspace{2mm} Cancelando términos y simplificando la expresión se obtiene:
$\quad \boldsymbol{med_{AB}:\  x-6y-9=0}$
 
\vspace{2mm} Es decir, cualquier pareja  $X=(x,y)$ que satisfaga la ecuación anterior $(x-6y-9=0)$ hará que la longitud entre $d(A, X)$ y $d(B, X)$ sea la misma (equidistantes).

\vspace{4mm}
\begin{center}
\rule{250pt}{0.1pt}

Más rápido: $\quad med_{AB} \, : \begin{cases} \ M=(3,-1) \\ \ \vec 	v_{med}=(6,1) \end{cases} \to \ \dfrac{x-3}{6}=\dfrac{y+1}{1} \ \to \ \boldsymbol{med_{AB}:\  x-6y-9=0}$
\end{center}

\end{miejemplo}

\vspace{5mm}

\subsection{Bisectriz}
\begin{tikzpicture}
	\fill [left color=red!50, right color=teal!50] (0,0) rectangle (3.5,.01);
	\fill [left color=teal!50, right color=blue!50] (3.5,0) rectangle (7.5,.01);
	\end{tikzpicture}
\vspace{0.5cm}

\begin{definition}[ Bisectriz]

La bisectriz de un ángulo es el lugar geométrico de los puntos del plano que equidistan de las rectas que forman el ángulo. 
	
\end{definition}

\begin{multicols}{2}
Gráficamente, una bisectriz divide un ángulo en dos partes iguales.  

$$bis_{\, r,s}\ = \ \{ \, X\in \mathbb R^2 \ \big/ \ d(X.r)=d(X,s) \, \}$$

Las bisectrices, rectas $b_1$ y $b_2$ forman con las rectas $r$ y $s$ la mitad del ángulo qye éstas forman.

$r:\ Ax+By+C=0;\ \ s:\ A'x+B'+y+C'=0 \quad \Rightarrow $

$\Rightarrow \ \  \dfrac{|A_1x+B_1y+C_1|}{\sqrt{A_1^2+B_1^2}}=\dfrac{|A_2x+B_2y+C_2|}{\sqrt{A_2^2+B_2^2}}$

\begin{figure}[H]
	\centering
	\includegraphics[width=0.5\textwidth]{img-lg/lg02.png}
\end{figure}
\end{multicols}


\begin{miejemplo}

Encontrar las ecuaciones de las bisectrices de los ángulos que determinan las rectas 

\textbf{\textit{3x - 4y + 5 = 0}}  $\ $ y  $\ $  \textbf{\textit{6x + 8y + 1 = 0}}.

\vspace{5mm}

Retomamos la definición de bisectriz  sustituimos
$\quad  \dfrac{|3x-4y+5|}{\sqrt{3^2+4^2}}=\dfrac{|6x+8y+1|}{\sqrt{6^2+8^2}}$
     
\vspace{2mm} Desarrollando tenemos:
$ \quad \dfrac{|3x-4y+5|}{5}=\dfrac{|6x+8y+1|}{10}$

\vspace{2mm} Es decir,
$\quad 10|3x-4y+5|=5|6x+8y+1|$
    
\vspace{2mm} De donde, $\quad 10(3x-4y+5)= \pm \big[ 5(6x+8y+1) \big]$ 

\vspace{2mm} Que conducen a las ecuaciones de los dos bisectrices,

$$\boldsymbol{ b_1:\ 16y-9=0 \ \qquad \ b_2:\ 12x+11=0}$$
\end{miejemplo}

\begin{miejercicio}

Calcula la ecuación de las bisectrices que forma la recta $r:\ 3x-4y+3=0$ con los ejes coordenados.

\rule{250pt}{0.1pt}

\vspace{2mm}$b_X:\ r \text{ con } OX \ \Rightarrow \ \dfrac{|3x-4y+3|}{\sqrt{3^2+(-4)^2}}=\dfrac{|y|}{\sqrt{0^2+1^2}} \ \to \ 3x-4y+3=\pm 5y \ \Rightarrow \ \begin{cases} \ b_{1X}:\ 3x-9y+3=0 \\ \ b_{2X}:\ 3x+y+3=0 \end{cases}$	

\vspace{5mm}$b_Y:\ r \text{ con } OY \ \Rightarrow \ \dfrac{|3x-4y+3|}{\sqrt{3^2+(-4)^2}}=\dfrac{|x|}{\sqrt{1^2+0^2}} \ \to \ 3x-4y+3=\pm 5x \ \Rightarrow \ \begin{cases} \ b_{1Y}:\ 2x+4y+3=0 \\ \ b_{2Y}:\ 8x-4y+3=0 \end{cases}$
\end{miejercicio}


\vspace{5mm}
\begin{cuadro-naranja}

\textbf{Bisectrices de dos rectas a partir de sus vectores directores}: este modo de encontrar las bisectrices será de mayor utilidad que el anterior en caso de querer encontrar las \emph{bisectrices interiores de un triángulo}, cuya intersección determinará el \emph{incentro}.

Usaremos \emph{vectores unitarios} $(\vec u_r,\ \ \vec u_s)$ para que el tamaño de los vectores directores $(\vec v_r,\ \ \vec v_s)$ no influya en la dirección de las bisectrices. (Bastaría con tomar vectores directores del mismo tamaño, por eso, la opción más sencilla es normalizarlos, hacerlos unitarios.

\begin{miejemplo}

Encuentra las bisectrices de las rectas $r:\, y=3x-2$ y $s:\, y=-2x+3$

\vspace{5mm} Es fácil encontrar que $\ P=r\cap s\equiv(1,1)$

\vspace{2mm} $\vec v_r=\ (m=3)\ =(1,3) \ \to \ \vec u_r=\dfrac{\vec v_r}{|\vec v_r|}=\dfrac {1}{\sqrt{10}} (1,3)$

\vspace{2mm} $\vec v_s=\ (m=-2)\ =(1,-2) \ \to \ \vec u_s=\dfrac{\vec v_s}{|\vec v_s|}=\dfrac {1}{\sqrt{5}} (1,-2)$

\vspace{2mm} Las direcciones de las bisectrices vendrán dadas por los vectores $\ \vec v_{b_1}=\vec u_r+\vec u_s \ \text{ y } \ \vec v_{b_2}=\vec u_r-\vec u_s\, . \ $ Así,

\begin{multicols}{2}
\vspace{2mm} $b_1:\ \begin{cases} \ P(1,1) \\ \ \vec v_{b_1}=\vec u_r+\vec u_s\approx (0.76,0.05) \end{cases}$

\vspace{3mm}$ \Rightarrow \ b_1:\ \dfrac{x-1}{0.76}=\dfrac{y-1}{0.05}$

\vspace{6mm} $b_2:\ \begin{cases} \ P(1,1) \\ \ \vec v_{b_2}=\vec u_r-\vec u_s\approx (-0.13,1.84) \end{cases} $

\vspace{3mm}$ \Rightarrow \ b_2:\ \dfrac{x-1}{-0.13}=\dfrac{y-1}{1.84}$

\begin{figure}[H]
	\centering
	\includegraphics[width=0.4\textwidth]{img-lg/lg03.png}
\end{figure}
\end{multicols}	
\end{miejemplo}
\vspace{1mm}
\end{cuadro-naranja}

\vspace{5mm}
\begin{figure}[H]
	\centering
	\includegraphics[width=1\textwidth]{img-lg/lg04b.png}
	\caption*{\textcolor{red}{\textbf{Bisectrices \emph{interiores} de un triángulo}}}
\end{figure}


\vspace{5mm}
\begin{miejercicio}

Encuentra las coordenadas del \emph{incentro} del triágulo de vértices $A(0,0),\ B(3,0) \text{ y } C(2,3)$

\rule{250pt}{0.1pt}

\vspace{6mm} $r_{AB}:\ y=0;\qquad r_{AC}:\ y=\dfrac 3 2 \, x;\qquad r_{BC}:\ y=-3x+9;\qquad \qquad A(0,0)$

\vspace{4mm} $b_{A}:\ \begin{cases} \ \vec v_{AB}=(3,0) &\to \vec u_{AB}=(1,0) \\ \ \vec v_{AC}=(2,3) &\to \vec u_{AC}=\dfrac{(2,3)}{\sqrt{13}}\approx (0.55,0.83) \end{cases}  \quad \Rightarrow \ $

\vspace{2mm} \hspace{3cm} $\Rightarrow \quad  \vec v_{b_A}=\vec u_{AB}+\vec u_{AC}=(1,0)+(0.55,0.83)=(1.55,0.83)$

\vspace{2mm} $b_{B}:\ \begin{cases} \ \vec v_{BA}=(-3,0) &\to \vec u_{BA}=(-1,0) \\ \ \vec v_{BC}=(-1,3) &\to \vec u_{BC}=\dfrac{(-1,3)}{\sqrt{10}}\approx (-0.32,0.95) \end{cases}  \quad \Rightarrow \ $

\vspace{2mm} \hspace{3cm} $\Rightarrow \quad  \vec v_{b_B}=\vec u_{BA}+\vec u_{BC}=(-1,0)+(-0.32,0.95)=(-1.32,0.95)$

\vspace{2mm} 
\begin{multicols}{2}
\vspace{4mm} $\vec v_{b_A} \ \to \ m_{b_A}=\ 0.54 \quad \ \Rightarrow \quad b_A:\ y=\ 0.54x$

\vspace{6mm} $\vec v_{b_B} \ \to \ m_{b_B}=-0.72 \quad \Rightarrow \quad b_B:\ y=-0.72x$

\vspace{8mm} $I\equiv b_A \cap b_B \ \to \ \begin{cases} \ y=0.54 x\\ \ y=-0.72x \end{cases} \ \to \ $

\vspace{4mm} $\to \ x=1.71;\ y=0.93 \ \Rightarrow \ \boldsymbol{I(1.71,0.93)}$
\begin{figure}[H]
	\centering
	\includegraphics[width=0.4\textwidth]{img-lg/lg05.png}
\end{figure}
\end{multicols}
	
\end{miejercicio}

%%%%%%%%%%%%%%%%%%%%%%%%%%%%%%%%%%%%%%%%%%%%%%%%%%%%%%%%%%%%%%%%%%%%%%%%
\vspace{5mm}

\section{Circunferencia}
\vspace{-5mm}
\begin{tikzpicture}
	\fill [left color=red!50, right color=teal!50] (0,0) rectangle (3.5,.1);
	\fill [left color=teal!50, right color=blue!50] (3.5,0) rectangle (7.5,.1);
	\end{tikzpicture}
\vspace{0.5cm}

\begin{definition}[ Circunferencia]

\begin{multicols}{2}
$\quad$ 

\textbf{\emph{Circunferencia}}: 

\vspace{4mm}es el lugar geométrico de los puntos del plano $X(x,y)$ cuya distancia a un punto fijo llamado centro, $\mathcal O(a,b)$, es una cantidad fija llamada radio,  $\mathcal R$.

\begin{figure}[H]
	\centering
	\includegraphics[width=.25\textwidth]{img-conicas/conicas20.png}
	\end{figure}
\end{multicols}
\vspace{-10mm}
$$\mathcal C \Big( \mathcal O(a,b)\, ,\  \mathcal R \Big)  \quad \equiv \quad \Big\{\,  X \in \mathbb R^2 \ \big/ \ d(X,\mathcal O)\ = \ \mathcal R \, \Big\}  \quad \leftrightarrow \quad \boldsymbol{ (x-a)^2 \ + \ (y-b)^2 \ = \ \mathcal R^2}$$
	
\end{definition}

\underline{Observaciones}:

Desarrollando: $\quad (x-a)^2+(y-b)^2=\mathcal R^2 \quad \longrightarrow \quad x^2+y^2-2ax-2by+a^2+b^2-\mathcal R^2=0$

\begin{itemize}

\item Se trata de un polinomio de segundo grado con dos variables, $x$ e $y$ y que no tiene término cruzado, en $x\cdot y$.
\item Los coeficientes de $x^2$ e $y^2$ son ambos $1$, por lo que lo prodremos escribir como:

$$ \boldsymbol{ x^2 \ + \ y^2\ + \ A\, x \ + \ B\, y \ + \ C \ = \ 0 }$$

con $\qquad -2a=A \ \to \ a=-A/2;\qquad -2b=B \ \to \ b=-B/2; \qquad $

$a^2+b^2-\mathcal R^2=C \ \Rightarrow \ \boldsymbol{ \mathcal R^2=\left( \dfrac A 2 \right)^2+\left( \dfrac B 2 \right)^2-C > 0 } \ (*)\ $ condición importante.

\item En estas condiciones, este polinomio de segundo grado representará una circunferencia centrada en $\mathcal O \left( -\dfrac A 2 ,- \dfrac B 2 \right)$ y de radio $\mathcal R=\sqrt{\left( \dfrac A 2 \right)^2+\left( \dfrac B 2 \right)^2-C }$	
\end{itemize}

Si la circunferencia está centrada en el origen su ecuación es $\ x^2+y^2=\mathcal R^2\ $ 
que se llama \emph{ecuación reducida de la circunferencia}.

\vspace{5mm} \underline{Propiedades de la circunferencia:}
\begin{itemize}
\item Dados dos puntos de un circunferencia, $P, \, Q \in \mathcal C$, entonces, la mediatriz del segmento $PQ$ pasa por el centro $\mathcal O$ de la circunferencia.	
\item Las tangentes en un punto a una circunferencia son perpendiculares al radio.
\item Por tres puntos (no alineados) pasa una única circunferencia, la que quiete su centro en el circuncentro del triángulo formado por los tres puntos.
\end{itemize}


\vspace{5mm} \begin{miejemplo}
 
Con los conocimientos adquiridos en temas anteriores, 	determina el lugar geométrico de los puntos del plano cuyas distancias a los puntos $A(0,0)$ y $B(6,3)$ están en razón de $2:1$. Interpreta el resultado obtenido.
 
 \vspace{5mm}  $\dfrac{d(X,A)}{d(X,B)}=2 \ \to \ \dfrac{\sqrt{x^2+y^2}}{\sqrt{(x-6)^2+(y-3)^2}}=2 \ \to \ x^2+y^2=4\, [(x-6)^2+(y-3)^2] \ \to \ $
 
\vspace{2mm} $\to \  x^2+y^2=2^2\, (x^2-12x+36 \, + \, y^2-6x+9) \ \to \ 3x^2+3y^2-48x-24y+180=0$
 
\vspace{2mm}  Veamos si se trata de una circunferencia, dividamos todo el polinomio por $3$ para conseguir que los coeficientes en $x^2$ e $y^2$ sean $1$, ya que no hay términos en $x\,y$:
 
\vspace{2mm}  $x^2+y^2-16x-8y+60=0\ $ 

\vspace{2mm} Para que realmente sea una circunferencia falta calcular el radio (si ello es posible): 

\vspace{2mm} $\quad \mathcal R=\sqrt{ \left(\dfrac{-16}{2} \right)^2+ \left(\dfrac{-8}{2} \right)^2 - 60} = \sqrt{64+16-60}=\sqrt{20}>0$, luego
 
\vspace{2mm}  Se trata de una circunferencia centrada en $(8,4)$ de radio $\sqrt{20}\, \mathrm{u}$.
 
 \end{miejemplo}

\vspace{5mm}

\begin{miejercicio}

Determina cuales de las siguientes ecuaciones determina una circunferencia y, en caso afirmativo, encuentra su centro y su radio.

\vspace{2mm} $a)\ \ x^2+y^2-4x+6=0;\qquad b)\ \ 3x^2+3y^2-12x+6y-12=0;\qquad c)\ \ x^2+y^2-4x-6y+13=0$

\rule{250pt}{0.1pt}

\vspace{2mm} $\triangleright \ \ a \quad $ Los términos en $x^2$ e $y^2$ son $1$ y no hay término cruzado $xy$. Calculemos el radio:

\vspace{2mm} $\mathcal R=\sqrt{(-4/2)^2+(0/2)^2-6}=\sqrt{-2} \ \to \ \nexists r$. No se trata de una circunferencia.


\vspace{4mm} $\triangleright \ \ b \quad $ Dividimos, primero, todo el polinomio por $3$:

\vspace{2mm} $x^2+y^2-4x+2y-4=0$ Como los términos en $x^2$ e $y^2$ son $1$ y no hay término cruzado $xy$, calculamos el radio:

\vspace{2mm} $\mathcal R=\sqrt{(-4/2)^2+(2/2)^2-(-4)}=\sqrt{9}=3$. Se trata de una circunferencia de centro $(-2,1)$ y radio $3$

\vspace{4mm} $\triangleright \ \ c \quad $ Como los términos en $x^2$ e $y^2$ son $1$ y no hay término cruzado $xy$, calculamos el radio:

\vspace{2mm} $\mathcal R=\sqrt{(-4/2)^2+(-6/2)^2-13}=0 \to \ \nexists r$, no se trata de una circunferencia.

\color{gris}
\vspace{4mm} Se podría resolver el problema completando cuadrados, para ello hay que fijarse en los términos en $x$ e $y$ que serán los dobles productos:	(Teorema 4.2, resolución de la ecuación de segundo grado) (ver al final del ejercicio)

\color{gris}
\vspace{2mm} $a)\quad (x-2)^2-4 \, + \, y^2 \, + 6 = 0 \ \to \ (x-2)^2+y^2=-2 <0\ \to \ $ No circunferencia



\color{gris}
\vspace{2mm} \textcolor{gris}{$b) \quad (x-2)^2-4\, + \, (y+1)^2-1\, -4 =0 \ \to \ (x-2)^2+(y+1)^2=9=3^2>0 \ $ Se trata de una circunferencia de centro $(2,-1)$ y radio $3\, \mathrm{u}$.}

\color{gris}
\vspace{2mm} $c)\quad (x-2)^2-4\, + \, (y-3)^2-9 \, + 13 = 0 \ \to \ (x-2)^2+(y-3)^2 \not > 0 \ $ No es una circunferencia. 

\color{gris}
\vspace{2mm} En este último caso, la suma de dos cantidades positivas (elevadas al cuadrado) es cero si y solamente si ambas son cero $\ \to \ (x-2)=0 \ \to \ x= 2 \ \ \wedge \ \ (y-3)=0 \ \to \ y=3 \quad \Rightarrow \quad P(2,3)\ $ La solución solo es un punto del plano.

\color{NavyBlue}
\begin{adjustwidth}{20pt}{20pt}
\vspace{6mm}\underline{Completar cuadrados}:

\vspace{2mm} $x^2+y^2-4x+5y+6=0 \ \to \ x^2-4x\ +y^2+5y \ +6=0 \ \to \ $

$\to \ x^2-2\cdot 2x + \textcolor{red}{ 2^2 - 2^2 }\ +y^2+2\cdot \frac 52 y + \textcolor{red}{(5/2)^2-(5/2)^2}\ +6=0  \ \to \ $

$\to \ (x-2)^2 \ -4 \ + \ (y+ 5/2)^2 \ - \ (5/2)^2\ + 6 = 0 \ \to \ $

$(x-2)^2+(y+5/2)^2\ -4+25/4+6 =0 \ \to \ (x-2)^2+(y+5/2)^2+33/4 = 0 \ \to $

$\to (x-2)^2+(y+5/2)^2=-33/4 < 0 \quad $, en este caso, no es una circunferencia.
\end{adjustwidth}
\color{black}
\end{miejercicio}

\vspace{0.5cm}

\subsection{Posiciones relativas de una recta $\boldsymbol r$ y una circunferencia $\boldsymbol{\mathcal C(\mathcal O,r)}$}
\begin{tikzpicture}
	\fill [left color=red!50, right color=teal!50] (0,0) rectangle (3.5,.01);
	\fill [left color=teal!50, right color=blue!50] (3.5,0) rectangle (7.5,.01);
	\end{tikzpicture}
\vspace{0.5cm}
\begin{figure}[H]
	\centering
	\includegraphics[width=1\textwidth]{img-conicas/conicas21.png}
	\end{figure}


$\checkmark \quad$ \underline{Algebráicamente}:

Resolver el sistema de ecuaciones formado por la recta y por la circunferencia,

\hspace{2cm} --- Si no hay solución $\quad \Rightarrow \quad r \text{ externa a } \ \mathcal C$ 

\hspace{2cm} --- Si la solución es única $\quad \Rightarrow \quad r \text{ tangente a } \ \mathcal C$

\hspace{2cm} --- Si hay dos soluciones $\quad \Rightarrow \quad r \text{ secante a } \ \mathcal C$

\vspace{4mm} $\checkmark \quad$ \underline{Geométricamente}:

Calculamos $d=d(\mathcal O, r)$, siendo $\mathcal O$ el centro de la circunferencia $\mathcal C$

\hspace{1cm}--- Si $\ d>\mathcal R \ \to \ r \text{ externa a } \ \mathcal C\, ; \ $ 

\hspace{1cm}--- Si $\ d=\mathcal R \ \to \ r \text{ tangente a } \ \mathcal C\, ; \ $ 

\hspace{1cm}--- Si $\ d<\mathcal R \ \to \ r \text{ secante a } \ \mathcal C\, ; \ $ 
$\quad$  Si $\ d=0 \ \to \ r \text{ secante a } \ \mathcal C\, ; \ $  en un diámetro.

\vspace{5mm}

\begin{miejemplo}

Determina la posición relativa de la circunferencia $\mathcal C:\, x^2+y^2-6x-4y-12=0$	 respecto de las rectas:

\vspace{2mm}$a)\ \ r_1:\, 3x-4y-26=0;\qquad b)\ \ r_2:\, 3x-4y-1=0;\qquad c)\ \ r_3:\, 3x-4y+49=0;$

\vspace{2mm}$d)\ \ r_4:\, x+2=0;\qquad \qquad \ \ \ e)\ \ r_5:\, y=4$

\vspace{6mm} $x^2+y^2-6x-4y-12=0 \quad \to \quad (x-3)^2-9+(y-2)^2-4-12=0 \ $

\vspace{2mm} \hspace{3cm} $\Rightarrow \quad (x-3)^2+(y-2)^2=5^2 \quad \Rightarrow \quad \mathcal C\, \Big( \, \mathcal O(3,2)\, ; \ \mathcal R=5\, \Big)$

\vspace{6mm} $\triangleright \ \ a)\quad d(\mathcal O,r_1)=\dfrac{|3\cdot 3-4\cdot 2-26|}{\sqrt{3^2+(-4)^2}}=\dfrac{25}{5}=5=\mathcal R \ \to \ $ tangentes.
 
\vspace{4mm} $\triangleright \ \ b)\quad d(\mathcal O,r_2)=\dfrac{|3\cdot 3-4\cdot 2-1|}{\sqrt{3^2+(-4)^2}}=0<\mathcal R \ \to \ $ secantes (diámetro).

\vspace{4mm} $\triangleright \ \ c)\quad d(\mathcal O,r_1)=\dfrac{|3\cdot 3-4\cdot 2+49|}{\sqrt{3^2+(-4)^2}}=\dfrac{50}{5}=10>\mathcal R \ \to \ $ externas.

\vspace{4mm} $\triangleright \ \ d)\quad d(\mathcal O,r_4)=\dfrac{|3+2|}{\sqrt{1^2+0^2}}=5=\mathcal R \ \to \ $ tangentes.

\vspace{4mm} $\triangleright \ \ f)\quad d(\mathcal O,r_r)=\dfrac{|2-4|}{\sqrt{0^2+1^2}}=2<\mathcal R \ \to \ $ secantes.

\vspace{6mm} Sería conveniente resolver el ejercicio algebraicamente.
\end{miejemplo}

\begin{miejercicio}

?`Para que valor de $k$, la recta $y=x+k$ es tangente a la crcinferencia $x*2+y^2=4$?	

\rule{250pt}{0.1pt}

\vspace{2mm} $r:\, x-y+k=0;\qquad x^2+y^2=2^2 \ \equiv \ \mathcal C \, \Big(\, \mathcal O(0,0)\, , \, \mathcal R=2\, \Big)$

\vspace{3mm} $d(\mathcal O,r)=\dfrac {|0-0+k|}{\sqrt{1^2+(-1)^2}}=\dfrac{|k|}{2} \ = \ \mathcal R=2 \quad \Rightarrow \quad k=\pm 2\sqrt{2}$
\end{miejercicio}


%\vspace{0.5cm}

\subsection{Potencia de una circunferencia}
\begin{tikzpicture}
	\fill [left color=red!50, right color=teal!50] (0,0) rectangle (3.5,.01);
	\fill [left color=teal!50, right color=blue!50] (3.5,0) rectangle (7.5,.01);
	\end{tikzpicture}
\vspace{0.5cm}


\begin{definition}[ Potencia]

Sea $P(x_0,y_0)$ un punto del plano y $\mathcal C \, \Big( \, \mathcal O(a,b)\, , \, \mathcal R \, \Big)$ una circunferencia.	 Calculamos la distancia entre $P$ y el centro $\mathcal O$ de la circunferencia y le llamamos $d=d(\mathcal O, P)$

\vspace{4mm} Se llama \emph{\textbf{potencia}}, $\boldsymbol{\mathcal P}$, de $P$ respecto de $\mathcal C$ a la cantidad $\boldsymbol{d^2-r^2}$

$$\boldsymbol{ \mathcal P \ =} \ d^2-r^2 \ \boldsymbol{ = \ (x_0-a)^2+(y_0-b)^2-r^2}$$

\vspace{4mm} Consiste en sustituir las coordenadas de $P$ en la ecuación (todo a la izquierda) de $\mathcal C$

\vspace{2mm} \textcolor{gris}{Evidentemente, }

\textcolor{gris}{$\ d=d(\mathcal O,P)= |\overrightarrow{OP} |=\sqrt{(x_0-a)^2+(y_0-b)^2} \ \to \ \mathcal P= d^2-r^2=(x_0-a)^2+(y_0-b)^2-r^2$}
\end{definition}
\underline{Observaciones}:

\begin{itemize}
\item Si $P$ es exterior a $\mathcal C \ \to \ (d>r) \ \to \ \mathcal P>0$
\item Si $P$ pertenece a $\ \mathcal C \ \to \ (d=r) \ \to \ \mathcal P=0$
\item Si $P$ es interior a $\mathcal C \ \to \ (d<r) \ \to \ \mathcal P<0$
\end{itemize}

\vspace{5mm}

\begin{definition}[ Eje radical]

Se llama \textbf{\emph{eje radical}} de dos circunferencia al lugar geométrico de los puntos del plano que tienen la misma potencia respecto de las dos circunferencias. Es una recta perpendicular al segmento que une sus centros.	
\end{definition}

\color{gris}

Por medio de la geometría analítica se puede demostrar que el eje radical es siempre una recta.

Sean $(a_1,b_1)$ y $(a_2,b_2)$ los centros de las circunferencias $\mathcal C_1$ y $\mathcal C_2$, respectivamente, de centros $\mathcal R_1$ y $\mathcal R_2$. Por definición de eje radical, un punto $(x,y)\in \mathbb R^2$ será del eje radical si:

$(x-a_1)^2+(y-b_1)^2-\mathcal R_1^2=(x-a_2)^2+(y-b_2)^2-\mathcal R_2^2 \ \to \ 
2(a_2-a_1)\, x \, + \, 2(b_2-b_1)\, y \, + \, (\mathcal R_2^2-\mathcal R_1^2+a_1^2+b_1^2-a_2^2-b_2^2)\ = \ 0$

Que es la ecuación de una recta $(Ax+By+C=0)$, pues no quedan términos de segundo grado $(x^2, \ y^2,\ x\cdot y)$.

El demostrar que es perpendicular al vector $\overrightarrow{\mathcal O_1 \mathcal O_2}$ es sencillo y se deja como ejercicio.

\color{black}

\vspace{5mm}

\begin{miejemplo}

Calcula la potencia del punto $P(-3,8)$ respecto de	las circunferencias:

\vspace{2mm} $\mathcal C_1:\ x^2+y^2-14x+13=0;\quad \mathcal C_2:\, \Big( \, \mathcal O(0,6)\, , \, \mathcal R=\sqrt{13} \, \Big); \quad \mathcal C_3:\  x^2+(y-6)^2=16$

E in I

\vspace{6mm} $\mathcal C_1:\ (x-7)^2-49+y^2+13=0 \ \to \ (x-7)^2+y^2=6^2$

\vspace{2mm} $\mathcal P_1= ((-3)-7)^2+8^2-6^2 = 100+64-36=128>0 \ \to \ P \ \text{ exterior a } \ \mathcal C_1$

\vspace{4mm} $\mathcal C_2:\ x^2+(y-6)^2-(\sqrt{13})^2=0 $

\vspace{2mm} $\mathcal P_2= (-3)^2+(8-6)^2-13 = 9+4-13=0 \ \to \ P \ \text{ pertenece a  } \ \mathcal C_1$

\vspace{4mm} $\mathcal C_3:\ x^2+(y-6)^2-4^2=0 $

\vspace{2mm} $\mathcal P_2= (-3)^2+(8-6)^2-16 =9+4-16=-3<0 \ \to \ P \ \text{ interior a } \ \mathcal C_1$
\end{miejemplo}

\begin{miejercicio}

Encuentra el eje radical de las circunferencias $\ \mathcal C_1:\ x*2+y^2-6x+4y-11=0\ $ y $\ \mathcal C_2:\ x^2+y^2+8x-2y-11=0\ $ y comprueba que se trata de una recta perpendicular al segmento que une sus centros.

\vspace{6mm} $P(x,y)\in \mathbb R^2 \ \to \ \begin{cases} \ \mathcal P_1(P,\mathcal C_1)=x^2+y^2-6x+4y-11 \\ \ \mathcal P_2(P,\mathcal C_2)=x^2+y^2+8x-2y-11 \end{cases} \quad \to \quad  \mathcal P_1=\mathcal P_2 \quad \Rightarrow$

\vspace{3mm}$\Rightarrow \quad 7x-3y+5=0 \ $ es una recta de vector director $\ \vec v=(3,7)$ \QED

\vspace{2mm} $\mathcal C_1:\ x^2+y^2-6x+4y-11=0 \ \to \ (x-3)^2-9+(y+2)^2-4-11=0 \ \to \ (x-3)^2+(y+2)^2=(\sqrt{24})^2 \ \Rightarrow \ \mathcal C_2 \, \Big( \, \mathcal O_1(3,-2)\, , \, \mathcal R_1=\sqrt{24}\, \Big)$

\vspace{2mm} $\mathcal C_2:\ x^2+y^2+8x-2y-11=0 \ \to \ (x+4)^2-16+(y-1)^2-1-11=0 \ \to \ (x+4)^2+(y-1)^2=(\sqrt{28})^2 \ \Rightarrow \ \mathcal C_2 \, \Big( \, \mathcal O_1(-4,1)\, , \, \mathcal R_1=\sqrt{28}\, \Big)$

\vspace{2mm} $\overrightarrow {\mathcal O_1 \mathcal O_2}=\mathcal O_2-\mathcal O_1=(-4,1)-(3,-2)=(-7,3)$

\vspace{2mm} $\vec v \cdot \overrightarrow {\mathcal O_1 \mathcal O_2} = (3,7)\cdot (-7,3)=3(7)+7(3)=-21+21=0 \ \Rightarrow \ \vec v \, \perp \,  \overrightarrow {\mathcal O_1 \mathcal O_2} \ \ $ Luego el eje radical es perpendicular al segmento que une los centros de las circunferencias que lo definen. \QED
	
\end{miejercicio}

\vspace{5mm}

\begin{myalertblock}{Circunferencia que pasa por tres puntos}

\vspace{1mm} Si $P,\, Q,\ \text{ y } \, R	\ $ son tres puntos no alineados, formarán un triángulo. Si encontramos su \emph{circuncentro} podremos trazar la circunferencia que pasa por los tres vértices del triángulo.

\vspace{4mm} $\checkmark \ \ $ \underline{Geométricamente}:

\vspace{2mm} $\qquad \qquad mtz_{\, PQ} \, \cap \, mtz_{\, PR} \ = \ \mathcal O\, ; \quad d(P,\mathcal O) \ = \ \mathcal R \quad \Rightarrow \quad \mathcal C \, \Big(\, \mathcal O \, , \,  \mathcal R\, \Big)$

\vspace{2mm} El centro lo encontraremos como la intersección de dos cualesquiera de las mediatrices del triángulo $PQR$. El radio no es más que la distancia del centro a cualquiera de los vértices.

\vspace{4mm} $\checkmark \ \ $ \underline{Algebraicamente}:

\vspace{2mm} $\qquad \qquad \mathcal C:\ x^2+y^2+Bx+Cy+D=0\ $ Imponiendo que $\ P,Q,R\, \in \ \mathcal C\, ,  $ tendremos un sistema de tres ecuaciones con tres incógnitas $A,B,C$ que determinarán $\mathcal C$

\vspace{1mm}

\end{myalertblock}

\vspace{5mm}

\begin{miejercicio}

Determina la ecuación de la circunferencia que pasa por los puntos $P(0,0),\ Q(10,0),\ R(18,12)$

\rule{250pt}{0.1pt}

\vspace{2mm} $\triangleright \ \ $ \underline{Geométricamente}:

\begin{small}
\vspace{2mm} $med_{\, PQ}=\begin{cases} \ M_{PQ}=\dfrac{P+Q}{2}=(5,0) \\ \perp \overrightarrow{PQ}=Q-P=(10,0)\ \to \ perp (1,0) \equiv (0,1) \end{cases} \to \ \dfrac{x-5}{0}=\dfrac{y}{1} \ \Rightarrow \ x=5$

\vspace{2mm} $med_{\, PR}=\begin{cases} \ M_{PR}=\dfrac{P+R}{2}=(9,6) \\ \perp \overrightarrow{PQ}=Q-P=(18,12)\ \to \ perp (3,2) \equiv (-2,3) \end{cases} \to \ \dfrac{x-9}{-2}=\dfrac{y-6}{3} \ \Rightarrow \ 3x+2y=39$
\end{small}

\vspace{2mm} $\ \mathcal O\,= \, mtz_{\, PQ} \, \cap \, mtz_{\, PR} \ = \begin{cases} \ x=5 \\ 3x+2y=39 \end{cases} \ \Rightarrow \ x=5,\ y=12 \ \Rightarrow \mathcal O(5,12)$

\vspace{2mm} $\mathcal R \, = \, d(\mathcal O, P) \, = \sqrt{(5-0 )^2+(12-0)^2}=\sqrt{169}=16\, \mathrm{u}$

\vspace{2mm} Finalmente, $\ \ \mathcal C:\ (x-5)^2+(y-12)^2=13^2$


\vspace{5mm} $\triangleright \ \ $ \underline{Algebraicamente}: $\quad \mathcal C:\ x^2+y^2+Ax+By+C=0$

\vspace{2mm} $P(0,0)\in \mathcal C \ \to \ 0^2+0^2+A\, 0+B\, 0+ C=0 \ \quad \ C=0$ 

\vspace{2mm} $Q(10,0)\in \mathcal C \ \to \ 10^2+0^2+A\, 10+B\, 0 + C=0 \ \quad \ 10A+C+100=0$

\vspace{2mm} $R(18,12)\in \mathcal C \ \to \ 18^2+12^2+A\, 18+B\, 12+ C=0 \ \quad \ 18A+12B+C+468=0$

\vspace{2mm} $\begin{cases} \ C=0 \\ \ 10A+C=-100 \\ \ 18A+12B+C=-468 \end{cases} \ \to \ C=0 \ \wedge \ A=-10 \ \wedge \ B=-24$

\vspace{2mm} Finalmente, $\quad \mathcal C:\ x^2+y^2-10x-24y=0 \ $
\begin{footnotesize} \textcolor{gris}{Comprobar que se trata de la misma circunferencia}. \end{footnotesize}


\end{miejercicio}

\begin{figure}[H]
	\centering
	\includegraphics[width=.7\textwidth]{img-conicas/conicas22.png}
	\caption*{\footnotesize{Ejercicio 12.6}}
	\end{figure}


%%%%
\begin{miejercicio}

Ecuación de la circunferencia $\mathcal C$, de centro $\mathcal O(2,3)$, que sea tangente a la recta $r:\, 4x-3y+6=0$

\rule{250pt}{0.1pt}

\vspace{2mm} $\mathcal R=d(\mathcal O,r)=\dfrac{|4(2)-3(3)+6|}{\sqrt{4^2+(-3)^2}}=1\, \mathrm{u}.\quad $ Luego, $\quad \mathcal C\, : \ (x-2)^2+(y-3)^2=1^2$
\end{miejercicio}


%%%%
\begin{miejercicio}

Los extremos de un diámetro de una circunferencia son $A(2,1)$ y $B(4,-5)$, encuentra su ecuación.

\rule{250pt}{0.1pt}

\vspace{2mm} $\mathcal O=M_{AB}=\dfrac{A+B}{2}=(3, -2)\, . \ \ $ Por otra parte, el diámetro $\phi=|\overrightarrow{AB}|=|B-A|=|(2,-6,)|=\sqrt{2^2+(-6)^2}=\sqrt{40}\, \mathrm{u} \ \to \ \mathcal R=\phi/2	=\sqrt{40}/2=\sqrt{10}\, \mathrm{u}\, , \ $ por lo que la circunferencia buscada tendrá por ecuación: $\ (x-3)^2+(y+2)^2=10$
\end{miejercicio}


%%%%
\begin{miejercicio}

Dada $\ \mathcal C:\ x^2+y^2-4x+10y+25=0$ determina la ecuación de una

\begin{enumerate}[a) ]
\item 	circunferencia concéntrica de radio $8$ u.
\item   circunferencia concentrica que pase por $P(2,2)$
\end{enumerate}

\rule{250pt}{0.1pt}

\vspace{2mm} $\mathcal C:\ x^2+y^2-4x+10y+25=0 \ \to \ (x-2)^2-4+(y+5)^2-25+25=0 \ \to \ $

\vspace{2mm} $\to \ (x-2)^2+(y+5)^2=2^2 \ \Rightarrow \ \mathcal C:\ \Big( \mathcal O(2,-5), \mathcal R=2 \Big)$

\vspace{4mm} $\triangleright \ \ a) \quad 	$ la circunferencia buscada tendrá centro en $(2,-5)$ y radio $8$ u, por lo que su ecuación será:
$\quad (x-2)^2+(y+y)^2=64$

\vspace{4mm} $\triangleright \ \ b) \quad 	$ la circunferencia buscada tendrá una ecuación de la forma:  $\ \ (x-2)^2+(y+5)^2=r^2\, , \  $ como queremos que pase por $P(2,2) \ \to \ \ (2-2)^2+(2+5)^2=r^2 \ \to \ 49=r^2 \ \Rightarrow \ r=7>0\, . \ $ La circunferencia buscada es: $\quad (x-2)^2+(y+5)^2=49$
\end{miejercicio}


%%%%
\begin{miejercicio}

Ecuación de la crcunferencia que pasa por el origen de coordenadas, tiene radio 4 u y su centro en la bisectriz del primer y tercer cuadrantes. \textcolor{gris}{$\ \ (y=x)$}

\rule{250pt}{0.1pt}

\vspace{2mm}	Si en centro de la circunferencia es $\mathcal O(a,b)$, al ser un punto de la bisectriz del primer y tercer cuadrantes, $y=x \ $ tendremos que $ \ a=b\, . \ $ Al tener radio $4$, su ecuación será: $\ \ (x-a)^2+(y-a)^2=16$

\vspace{2mm} Si exigimos ahora que pase por el origen de coordenadas $(0,0)\, , \  $ al sustituir en la ecuación tendremos $ \ (0-a)^2+(0-a)^2=16 \ \to \ 2a^2=16 \ \to \  a^2=8 \ \to \ a=\pm 2\sqrt{2}$

\vspace{2mm} Hay dos soluciones: $\ \ \mathcal C_1:\ (x-2\sqrt{2})^2+(y-2\sqrt{2})^2=16 \ \ \wedge \ \ \mathcal C_2:\ (x+2\sqrt{2})^2+(y+2\sqrt{2})^2=16 $
\end{miejercicio}



%%%%
\begin{miejercicio}

Ecuación de la recta tangente a $x^2+y^2=25$ que pasa por $A(0,6)$.

\vspace{2mm} Repetir el ejercicio si ahora $A(3,4)$

\rule{250pt}{0.1pt}


\vspace{2mm} $r$ tangente a $\mathcal C$ si $d(r,\mathcal O)=\mathcal R$, nuestra circunferencia tiene centro en el origen $\mathcal O(0,0)$ y radio $\mathcal R=5\, \mathrm{u}$

\vspace{2mm} En el primer caso $A\notin \mathcal C$, en el segundo $A\in \mathcal C$, pero procederemos del mismo modo en los dos casos.	

\vspace{4mm} $\triangleright \ \ $ Si $\ A(0,6) \quad \to \quad y-6=mx \quad \to \quad r:\ mx-y+6=0 \ \Rightarrow \ d(\mathcal O,r)=\dfrac{|6|}{\sqrt{m^2+1}}=5=\mathcal R \quad \to \quad 6=5\sqrt{m^2+1} \ \to \ 36=25(m^2+1) \ \to m^2=11/25 \ \to \ m=\pm \dfrac{\sqrt{11}}{5}\, , \ $ hay dos soluciones.

\vspace{4mm} $\triangleright \ \ $ Si $\ A(3,4) \quad \to \quad y-4=m(x-3) \quad \to \quad r:\ mx-y+3-4m=0 \quad \Rightarrow \quad $

$\Rightarrow \quad  d(\mathcal O,r)=\dfrac{|3-4m|}{\sqrt{m^2+1}}=5 \quad \to \quad (3m-4)^2=5^2(m^2-1) \ \cdots \ \ \to m=\dfrac{-4(1\pm \sqrt{2})}{3}\, , \ $ también hay dos soluciones.

\end{miejercicio}

%%%%
\begin{miejercicio}

Circunferencia que pasa por $A(2,1)$ y $B(4,-3)$ y tiene su centro en la recta $r:\ x+y-5=0$

\rule{250pt}{0.1pt}

\begin{multicols}{2}
El centro de la circunferencia lo encontraremos como intersección de la recta $r$ con la mediatriz del segmento $AB$. El radio no será más que la distancia del centro a uno cualquiera de los puntos $A$ y $B$ por donde pasa la circunferencia.

\vspace{2mm} $mtz_{AB}:\ \begin{cases} \ M_{AB} \\ \ \perp \overrightarrow{AB} \end{cases} \ \cdots \cdots \to \ \ x-2y=5$


\begin{figure}[H]
	\centering
	\includegraphics[width=.5\textwidth]{img-conicas/conicas23.png}
	\end{figure}	
\end{multicols}
 $\mathcal O = r \cap mtz_{AB}=\begin{cases} \ x+y=5 \\ \ x-2y=5 \end{cases} \ \to \ \mathcal O(5,0)$

\vspace{2mm} $\mathcal R=d(\mathcal O,A)=\cdots =	\sqrt{10} \quad \Rightarrow \quad \mathcal C:\ (x-5)^2+y^2=10$	
\end{miejercicio}


%%%%
\begin{miejercicio}

Recta tangente a la circunferencia $x^2+y^2=16$ que sea paralela a la recta $r:\ 2x-y+3=0$

\rule{250pt}{0.1pt}

\vspace{2mm}$\mathcal C \big( (0,0)\, ; \, 4\Big);\qquad r:\ y=2x+3 \ \to \ m_r=2$ La recta buscada, por ser paralela a $r$ tendrá por ecuación $s:\ y=2x+k$, en forma general, $2x-y+k=0$

\vspace{2mm} $s$ tangente a $ \mathcal C \ \leftrightarrow \ d(\mathcal O, s)=\mathcal R \quad \to \quad \dfrac{ |2(0)-(0)+k| }{ \sqrt{ 2^2+(-1)^2 } }=4 \  \to \ \dfrac{|k|}{\sqrt{5}}=4 \ \Rightarrow \ k=\pm 4\sqrt{5}\ $ hay dos soluciones.	
\end{miejercicio}


%%%%%%%%%%%%%%%%%%%%%%%%%%%%%%%%%%%%%%%%%%%%%%%%%%%%%%%%%%%%%%%%%%%%%%%%

\vspace{5mm}

\subsection{Cónicas (en el próximo capítulo)}
\begin{tikzpicture}
	\fill [left color=red!50, right color=teal!50] (0,0) rectangle (3.5,.01);
	\fill [left color=teal!50, right color=blue!50] (3.5,0) rectangle (7.5,.01);
	\end{tikzpicture}
\vspace{0.5cm}

La importancia y la extensión de las cónicas es lo suficientemente importante para dedicarles el próximo capítulo completo.

\vspace{1cm}



\section{Ejercicios}
\begin{tikzpicture}
	\fill [left color=red!50, right color=teal!50] (0,0) rectangle (3.5,.01);
	\fill [left color=teal!50, right color=blue!50] (3.5,0) rectangle (7.5,.01);
	\end{tikzpicture}
\vspace{0.5cm}

%%%
\begin{miejercicio}

Determina el lugar geométrico de los puntos del plano que distan 2 unidades de la recta $r:\, y=1$.

\rule{250pt}{0.5pt}

\vspace{2mm} $X(x,y) \ \big/ \ d \big( \ X(x,y)\, , \, r:\, y-1=0 \ \big) \ = \ 2 \quad \Rightarrow \ \left| \dfrac{y-1}{2} \right| = 2  \quad \to \quad \begin{cases} \ y=\ 3 \\ \ y=-1 \end{cases}$

\vspace{2mm} Se trata de dos rectas paralelas a $r$ a distancia de 2 unidades cada una de ellas.
	
\end{miejercicio}

%%%
\begin{miejercicio}

Determina el lugar geométrico de los puntos del plano que distan del eje de abcisas el triple que del eje de ordenadas.

\rule{250pt}{0.5pt}

\vspace{2mm} 
\begin{multicols}{2}
$d\,(X,\, eje\ X) \ =\ 3\ d\, (X,\, eje\ Y) \quad \to $

\vspace{5mm} $\quad \left| \dfrac y 1 \right| \ = \ \left| \dfrac x 1 \right| \ \Rightarrow \ \begin{cases} \ y=\ 3x \\ \ y=-3x \end{cases} $

\vspace{5mm} Son dos rectas que pasan por el origen y de pendientes $3$ y $-3$

\begin{figure}[H]
	\centering
	\includegraphics[width=0.2\textwidth]{img-lg/lg06.png}
\end{figure}
\end{multicols}
	
\end{miejercicio}

%%%
\begin{miejercicio}

Determina el lugar geométrico de los puntos del plano que equidistan de la bisectriz del primer y tercer cuadrantes (recta $y=x$) y del punto $A(2,0)$.

\rule{250pt}{0.5pt}

\vspace{2mm} $d\, (X.A) \ = \ d\, (X,r) \quad \to \quad \left| \dfrac{1\, x - 1\, y - 0}{\sqrt{1^2+(-1)^2}} \right|\ = \ \sqrt{(x-2)^2 +(y-0)^2}	\quad \to \quad $

\vspace{2mm} $\to \quad x-y=\pm \sqrt{2}\, \sqrt{(x-2)^2 +(y-0)^2}  \quad \to \quad (x-y)*2 =\pm 2\, [\, (x-2)^2 +y^2 \,] \quad \Rightarrow $

\vspace{2mm} $ \Rightarrow \quad x^2+y^2-2xy-8x+8=0 $
\end{miejercicio}

%%%
\begin{miejercicio}

Determina el lugar geométrico de los puntos del plano tales que las distancias a los puntos $A(-3,1)$ y $B(1,0)$ están en razón de $1:2$.

\rule{250pt}{0.5pt}

\vspace{2mm} $\dfrac{d(X,A)}{d(X,B)} \ = \ \dfrac 1 2 \quad \to \cdots \cdots \quad \to 3x^2+3y^2-26x-14y+30=0$
	
\end{miejercicio}

%%%
\begin{miejercicio}

Determina el lugar geométrico de los puntos del plano que equidistan del puntos $A(-1,4)$ (radio) una cantidad $r=5\, \mathrm{u}$ (radio). (Circunferencia).

\rule{250pt}{0.5pt}

\vspace{2mm} $d(X,A)=5 \quad \to \quad \sqrt{(x+1)^2+(y-4)^2}=5 \quad \Rightarrow \quad (x+1)^2+(y-4)^2=25 \quad \Rightarrow $

\vspace{2mm} $\Rightarrow  \quad \ x^2+y^2+2x-8y-8=0$
	
\end{miejercicio}

%%%
\begin{miejercicio}

Determina el lugar geométrico de los puntos del plano cuya suma de distancias a dos puntos fijos $F(-2,0)$ y $F'(2,0)$ sea igual a $5\, \mathrm{u}$ (hipérbola).

\rule{250pt}{0.5pt}

\vspace{2mm} $\sqrt{(x+2)^2+y^2}+\sqrt{(x-2)^2+y^2}=5 \ \to \ 
\sqrt{(x+2)^2+y^2}=5-\sqrt{(x-2)^2+y^2}\ $ elevando al cuadrado,

\vspace{2mm} $(x+2)^2+y^2=25+(x-2)^2+y^2-10\, \sqrt{(x-2)^2+y^2}\ $ aislando la raiz y elevando al cuadrado

\vspace{2mm} $(8x-25)^2=\big(-10\, \sqrt{(x-2)^2+y^2}\, \big)^2 \quad \to \quad
 36x^2+100y^2=225$
 
 \vspace{2mm} Ecuación que arreglando un poco da $\ \dfrac{4}{25}	x^2 + \dfrac 4 9 y^2 = 1 \quad \Rightarrow \quad \dfrac{x^2}{(5/2)^2}+\dfrac{y^2}{(3/3)^2}=1$
 
\end{miejercicio}

%%%
\begin{miejercicio}

Determina el lugar geométrico de los puntos del plano cuya resta de distancias a dos puntos fijos $F(-2,0)$ y $F'(2,0)$ sea igual a $5\, \mathrm{u}$ (elipse).

\rule{250pt}{0.5pt}

\vspace{2mm} $\sqrt{(x+2)^2+y^2}-\sqrt{(x-2)^2+y^2}=5 \ \to \ 
\sqrt{(x+2)^2+y^2}=5+\sqrt{(x-2)^2+y^2}\to  \  
\cdots \cdots \ \to $

\vspace{2mm} $\to \quad 44x^4-100y^2=275 \quad \Rightarrow \quad \dfrac{x^2}{(5/2)^2}-\dfrac{y^2}{(\sqrt{40/11})^2}=1$
	
\end{miejercicio}

%%%
\begin{miejercicio}

Determina el lugar geométrico de los puntos del plano tal que la diferencia de los cuadrados de las distancias a los puntos $A(0,0)$ y $B(3,1)$ es de $7\, \mathrm{u}$.

\rule{250pt}{0.5pt}

\vspace{2mm} $\big(d(X,A)\big)^2-\big(d(X,B)\big)^2=7 \quad \to \quad x^2+y^2-[\, (x-3)^2+(y-1)^2\, ]=7 \quad \to \quad 6x-2y+17=0$
	
\end{miejercicio}

\vspace{1cm}

\begin{mipropuesto}

Encuentra el incentro del triángulo de vértices $\ A(2,3),\ B(-1,2),\ C(3,0)$ y determina el radio de la circunferencia inscrita.	
\end{mipropuesto}

\vspace{-6mm}
\begin{flushright}
\begin{footnotesize} \textcolor{gris}{\rotatebox{180}{ $\textcolor{gris}{ r=d(I,r_{AB}) } \qquad \qquad I(\sqrt{2,2\sqrt{2}-1)};\quad r=\sqrt{10}-\sqrt{5}\, \mathrm{u}$ }}	\end{footnotesize}
\end{flushright}

\begin{mipropuesto}

Dado el triángulo de vértices $\ A(2,1),\ B(1,-2),\ C(-1,3)$, 

\begin{enumerate}[a) ]
\item Calcula $P$, intersección de la bisectriz interior de vértice $C$ con el lado opuesto $AB$.
\item Demuestra que $\ \dfrac{PA}{PB}=\dfrac{CA}{CB}$.	
\end{enumerate}
	
\end{mipropuesto}

\vspace{-6mm}
\begin{flushright}
\begin{footnotesize} \textcolor{gris}{\rotatebox{180}{ $P\left( \dfrac{45-\sqrt{377}}{16},\dfrac{55-3\sqrt{377}}{16} \right) \, ; \quad  \text{ sí se cumple, } \ \dfrac{PA}{PB}=\dfrac{CA}{CB}=\sqrt{ \dfrac{13}{29} }$ }}	\end{footnotesize}
\end{flushright}

%%
\begin{mipropuesto}

Averigua cuáles de las siguientes expresiones corresponden a circunferencias y, en ellas, halla su centro y su radio:

\hspace{1cm} \vspace{2mm} $a)\ \ m x^2 + y^2 - 8x + 2y + 10 = 0\ ; \qquad   \ b) \ \ x^2-y^2 + 2x + 3y -5 = 0$


\hspace{1cm} \vspace{2mm} $c)\ \  x^2 + y^2 + xy -x + 4y - 8 = 0 \ ; \qquad d)\ \ 2x^2 + 2y^2 -16x + 24 = 0$
	
\end{mipropuesto}

\vspace{-6mm}
\begin{flushright}
\begin{footnotesize} \textcolor{gris}{\rotatebox{180}{ $a)\ \ \mathcal C \Big( (4,-1);\sqrt{7}\Big);\quad b)\ $ No, coeftes. $x^2$ e $y^2$ distintos; $\quad c)\ $ No, hay tmno. en $xy$; $\quad d)\ \ \mathcal C \Big( (4,0);2 \Big)$ }}	\end{footnotesize}
\end{flushright}

%%
\begin{mipropuesto}

Uno de los diámetros de una circunferencia tiene por extremos $A(3, -2)$ y $B(7, 0)$. Halla la ecuación de la circunferencia. 
	
\end{mipropuesto}

\vspace{-6mm}
\begin{flushright}
\begin{footnotesize} \textcolor{gris}{\rotatebox{180}{ $(x-5)^2+(y+1)^2=5$ }}	\end{footnotesize}
\end{flushright}

%%
\begin{mipropuesto}

Determina la ecuación de la circunferencia que pasa por $A(2, - 4),\  B(8, -10) \text{ y } C(4, -8)$ 
	
\end{mipropuesto}

\vspace{-6mm}
\begin{flushright}
\begin{footnotesize} \textcolor{gris}{\rotatebox{180}{ $(x-9)^2+(y+3)^2=50$ }}	\end{footnotesize}
\end{flushright}

%%
\begin{mipropuesto}

Estudia la posición relativa de la circunferencia de ecuación $x^2 + y^2 - 6x - 4y + 9 = 0$ respecto de cada una de las siguientes rectas: $\ r_1: \ x + y –-1 = 0 ;\quad r_2:\  3x - 4y + 9 = 0$
	
\end{mipropuesto}

\vspace{-6mm}
\begin{flushright}
\begin{footnotesize} \textcolor{gris}{\rotatebox{180}{ Exterior y tangente. }}	\end{footnotesize}
\end{flushright}

%%
\begin{mipropuesto}

Halla el eje radical de los siguientes pares de circunferencias: 

\vspace{2mm} \hspace{1cm} $a)\  \ x^2 + y2 = 4\ \ $ y $\ \ x^2 + ( y - 1)^2 = 9$

\vspace{1mm} \hspace{1cm} $b)\  \ (x - 3)^2 + y^2 = 5 \ \ $  y $\ \ (x -7)^2 + y^2 = 9$

\vspace{1mm} \hspace{1cm} $c)\ \ x^2 + ( y - 3)^2 = 2\ \ $ y $\ \ (x -5)^2 + y^2 = 1$ 
	
\end{mipropuesto}

\vspace{-6mm}
\begin{flushright}
\begin{footnotesize} \textcolor{gris}{\rotatebox{180}{ $a)\ \ y=-2;\qquad b)\ \ x=9/2;\qquad c)\ \ 5x-3y-7=0$ }}	\end{footnotesize}
\end{flushright}

%%
\begin{mipropuesto}

De la circunferencia $\mathcal C$ se sabe que tiene su centro en la recta $r:\ x - 3y = 0$ y pasa por los puntos $A(1, 4)$ y $B(3, 6)$. Obtén la ecuación de $\mathcal C$.
	
\end{mipropuesto}

\vspace{-6mm}
\begin{flushright}
\begin{footnotesize} \textcolor{gris}{\rotatebox{180}{ $\mathcal O(3t,t)\;\quad  t \ \big/  \ d(\mathcal O,A)=d(\mathcal O,B) \quad  \to \quad \mathcal C:\ (x-3)^2+(y-1)^2=25$ }}	\end{footnotesize}
\end{flushright}









\vspace{20mm}
\begin{adjustwidth}{50pt}{250pt}
\begin{cuadro-naranja}
\textbf{\huge{Problemas $\boldsymbol{+}$}}\normalsize{$\, $}
\end{cuadro-naranja}	
\end{adjustwidth}

\vspace{5mm}
\begin{enumerate}[\textbf{P$\boldsymbol +$} 1. ]


%%%%
\item	Halla el lugar geométrico de los puntos del plano que equidistan de los siguientes pares de rectas:

\vspace{3mm} $a)\ \ \begin{cases} \ r_1:\, 3x-4y-26&=0 \\ \ r_2:\, 12x+5y+1&=0 \end{cases} \qquad \qquad b)\ \ \begin{cases} \ s_1:\,  2x\ -7y\ -9&=0 \\ \ s_2:\, 4x-14y+11&=0 \end{cases}$


\vspace{-4mm}
\begin{flushright}
\begin{footnotesize} \textcolor{gris}{\rotatebox{180}{ $b)\ \ 8x-28y-7=0\quad $ Como $s_1\parallel s_2 \to \ $ se obtiene una recta paralela a ambas que equidista de ellas. }}	\end{footnotesize}
\end{flushright}

\vspace{-6mm}
\begin{flushright}
\begin{footnotesize} \textcolor{gris}{\rotatebox{180}{ $a)\ \ \begin{cases} \ 3x-11y+49&=0 \\ \ 11x-3y-37&=0 \end{cases} \quad $ Son las birectrices de $r_1$ y $r_2$. }}	\end{footnotesize}
\end{flushright}

%%%%
\item	Dados los puntos $A(2,1)$ y $B(6,4)$, determina el lugar geométrico de los puntos $P$ del plano tales que la superficie del triángulo $ABP$ sea de 10 unidades de área.

\vspace{-6mm}
\begin{flushright}
\begin{footnotesize} \textcolor{gris}{\rotatebox{180}{ La solución son dos rectas paralelas: $\quad 3x-4y-22=0;\qquad 3x-4y+18=0$ }}	\end{footnotesize}
\end{flushright}
\vspace{-6mm}
\begin{flushright}
\begin{footnotesize} \textcolor{gris}{\rotatebox{180}{ $AB$ base del triángulo, de longitud 5, luego la altura, $d(P,r_{AB})=$ 4.}}	\end{footnotesize}
\end{flushright}

\item Estudia la posición relativa de la recta $y = x + k$ y la circunferencia $x^2 + y^2 = 1$ en función del parámetro $k$.

\vspace{-6mm}
\begin{flushright}
\begin{footnotesize} \textcolor{gris}{\rotatebox{180}{$k=\pm \sqrt{2} \ \to \ \text{ tangente}; \qquad k\in ]-\sqrt{2},\sqrt{2}[ \ \to \ \text{ secante}$; $\qquad$ otros casos: exterior.  }}	\end{footnotesize}
\end{flushright}

\item Considera las circunferencias $\ \mathcal C_1:\  (x - 1)^2 + ( y + 1)^2 = 2 \  \text{ y } \ \mathcal C_2:\  (x - 3)^2 + ( y + 3)^2 = 10$.

a) Comprueba que ambas circunferencias son secantes y calcula sus puntos de corte, $A$ y $B$. \textcolor{gris}{Resuelve el sistema formado por las dos ecuaciones de las circunferencias.}

b) Halla las potencias de los puntos $A$ y $B$ a las circunferencias $\mathcal C_1$ y $\mathcal C_2$.

c) A la vista del resultado obtenido en el apartado anterior, ?`qué podrías decir del eje radical de ambas circunferencias?

\vspace{-6mm}
\begin{flushright}
\begin{tiny} \textcolor{gris}{\rotatebox{180}{ $a)\ \ A(0,-2),\ B(2,0);\qquad b)\ \ \mathcal P(A,\mathcal C_1)=\mathcal P(A,\mathcal C_2)=0 \ \wedge \ \mathcal P(b,\mathcal C_1)=\mathcal P(B,\mathcal C_2)= 0;\ \ \text{secantes};\qquad c)\ \ \text{ por ello, eje radical:} \ r_{AB}$ }}	\end{tiny}
\end{flushright}

\item  Halla la ecuación de la circunferencia inscrita en el cuadrado de vértices $A(-3, 3)$, $B(-1, 3)$, $C(-1, 1)$ y $D(-3, 1)$.

\vspace{-6mm}
\begin{flushright}
\begin{footnotesize} \textcolor{gris}{\rotatebox{180}{ $\mathcal O=r_{AC}\cap r_{BD};\ \ \mathcal R=\frac 12 d(A,B)  \quad \to \quad \mathcal C:\ (x+2)^2+(y-2)^2=1$}}	\end{footnotesize}
\end{flushright}

\end{enumerate}

\begin{figure}[H]
	\centering
	\includegraphics[width=0.5\textwidth]{img-lg/lg08.png}
\end{figure}




\newpage
\vspace{1cm}
\section{Resumen del tema}
\begin{tikzpicture}
	\fill [left color=red!50, right color=teal!50] (0,0) rectangle (3.5,.1);
	\fill [left color=teal!50, right color=blue!50] (3.5,0) rectangle (7.5,.1);
	\end{tikzpicture}
\vspace{0.5cm}

\begin{myblock}{Resumen ``lugares geométricos''}

$$\text{Mediatriz segmento } \overline{AB} \quad \equiv \quad \{ \, X(x,y)\in \mathbb R^2\ \big/ \ d(X,A) \ = \ d(X,B) \, \}$$

\vspace{5mm} $$\text{Med}_{AB} \quad \equiv \quad \begin{cases} \ M\ = \ \dfrac{A+B}{2} \\  \ \perp \ \overrightarrow{AB} \end{cases}$$ 

\vspace{5mm}
\begin{center} \rule{200pt}{0.1pt}  \end{center}


\vspace{5mm}$$\text{Bisectriz}(r,s) \quad \equiv \quad  \{ \, X\in \mathbb R^2 \ \big/ \ d(X.r)=d(X,s) \, \}$$

\vspace{5mm} $$A,\, B,\, C,\ \text{ triángulo } , \quad \text{Bisectriz interna } A \quad \equiv \quad \begin{cases} \ A \\ \ \vec v_a=\vec u_{AB}+\vec u_{AC} \end{cases} \quad $$

$$\text{ con } \quad \vec u_{AB}=\dfrac{\overrightarrow{AB}}{|\overrightarrow{AB}|} \ \ \ \wedge \ \ \  \vec u_{AC}=\dfrac{\overrightarrow{AC}}{|\overrightarrow{AC}|} $$

	
\end{myblock}



\begin{comment}

%%%%%%%%%%%%%%%%%%%%%%%%%%%%%%%%%%%. SECCIONES
\begin{figure}[H]
	\centering
	\includegraphics[width=0.6\textwidth]{ruta/imagen.png}
\end{figure}


\chapter{texto}
\begin{tikzpicture}
	\fill [left color=red!50, right color=teal!50] (0,0) rectangle (6.5,.2);
	\fill [left color=teal!50, right color=blue!50] (6.5,0) rectangle (11.5,.2);
	\end{tikzpicture}

\vspace{1cm}
\section{texto}
\begin{tikzpicture}
	\fill [left color=red!50, right color=teal!50] (0,0) rectangle (3.5,.1);
	\fill [left color=teal!50, right color=blue!50] (3.5,0) rectangle (7.5,.1);
	\end{tikzpicture}
\vspace{0.5cm}

\subsection{texto}
\begin{tikzpicture}
	\fill [left color=red!50, right color=teal!50] (0,0) rectangle (3.5,.01);
	\fill [left color=teal!50, right color=blue!50] (3.5,0) rectangle (7.5,.01);
	\end{tikzpicture}
\vspace{0.5cm}


%%%%%%%%%%%%%%%%%%%%%%%%%%%%%%%%%%%. \begin{ ------>. 
detsacado;  cuadro-naranja;  cuadro-gris;  miejercicio (solución extensa);  mipropuesto (solución corta y fuera del cuadro)

%%%%%%%%%%%%%%%%%%%%%%%%%%%%%%%%%%%. CURIOSIDAD
\vspace{1cm}
\color{ForestGreen!80}
\rule{250pt}{0.2pt}
Texto
\vspace{-8mm}
\begin{flushright}
\rule{250pt}{0.2pt}		
\end{flushright}	
\color{black}
\end{comment}